\documentclass[12pt]{amsart}
\usepackage{amsmath,amsthm,amssymb}
\usepackage{mathtools}
\usepackage[margin=1.15in]{geometry}
\newtheorem{lem}{Lemma}
\newtheorem*{lemA}{Lemma (Hrushovski--Loeser, 8.5.2)}
\theoremstyle{remark}
\newtheorem*{rem}{Remark}
\newcommand{\Z}{\mathbb Z}
\newcommand{\Q}{\mathbb Q}
\newcommand{\N}{\mathbb N}
\begin{document}
\title{A direct proof of Lemma 8.5.2 of Hrushovski--Loeser}
\date{November 20, 2026. Working note for Section 8; not part of the manuscript.}
\maketitle

Let $R = \Z[U, U^{-1}]$ and let $R[T]_\dag$ be the localization of $R[T]$ at the multiplicative set generated by the polynomials $1 - U^a T^b$, $a \in \Z$, $b \geq 1$, viewed inside $R[[T]]$. For $h = P/Q$ in $R[T]_\dag$ put $\deg h = \deg P - \deg Q$. If $\deg h \leq 0$ then $\lim_{T \to \infty} h$ is defined: it is $0$ if $\deg h < 0$, and the quotient of the leading coefficients of $P$ and $Q$ if $\deg h = 0$. Equivalently, it is the constant term of the expansion of $h$ in powers of $T^{-1}$; so the $h$ with $\deg h \leq 0$ form a subgroup on which $\lim$ is additive.

The value group $\Gamma$ is a divisible ordered abelian group whose definable points form $\Q$, the value of $t$ being $1$. A definable subset of $\Gamma^n$ is a Boolean combination of sets $g > 0$, $g = 0$ with $g$ affine with rational coefficients. We write $\chi$ for the Euler characteristic of definable sets.

\begin{lemA}
Let $\Delta \subseteq \Gamma^n$ be bounded and definable, and let $\ell : \Delta \to \Gamma$ be, on each piece of a finite definable partition of $\Delta$, of the form $x \mapsto \sum_i a_i x_i + b$ with $a_i, b \in \Z$. Put
\[
Z(\Delta, \ell)(T) = \sum_{m \geq 1} \ \sum_{\gamma \in \Delta \cap (m^{-1}\Z)^n} U^{-m \ell(\gamma)}\, T^m .
\]
Then $Z(\Delta, \ell) \in R[T]_\dag$, $\deg Z(\Delta,\ell) \leq 0$ and $\lim_{T \to \infty} Z(\Delta, \ell) = - \chi(\Delta)$.
\end{lemA}

An \emph{open simplex} of dimension $k$ is a set $\{\sum_{i=0}^k t_i v_i : t_i > 0, \ \sum_i t_i = 1\}$ with $v_0, \ldots, v_k$ affinely independent; here the $v_i$ are always in $\Q^n$ and the $t_i$ are in $\Gamma$ (recall that $1 \in \Gamma$ is the value of $t$). An open simplex of dimension $k$ has Euler characteristic $(-1)^k$.

\begin{lem}[Prism]\label{prism}
Let $\sigma \subseteq \Gamma^{n}$ be an open simplex with vertices $v_0, \ldots, v_k$ and let $f, g$ be rational affine functions with $f < g$ on $\sigma$. Then $\{(x, y) : x \in \sigma, \ f(x) < y < g(x)\}$ is a finite disjoint union of open simplices with rational vertices.
\end{lem}
\begin{proof}
Put $d_i = g(v_i) - f(v_i)$. Since $g - f$ is affine and positive on $\sigma$, we have $d_i \geq 0$ for all $i$, and not all $d_i$ vanish. For $-1 \leq j \leq k$ let $h_j$ be the affine function on $\sigma$ with $h_j(v_i) = f(v_i)$ for $i \leq j$ and $h_j(v_i) = g(v_i)$ for $i > j$. In barycentric coordinates, for $x = \sum_i t_i v_i$ we have $h_j(x) = f(x) + \sum_{i > j} t_i d_i$. Hence
\[
g = h_{-1} \geq h_0 \geq \cdots \geq h_k = f, \qquad h_{j-1} - h_j = t_j d_j ,
\]
and on $\sigma$ we have $h_{j-1} > h_j$ if $d_j > 0$ and $h_{j-1} = h_j$ if $d_j = 0$. So the set in question is the disjoint union of the bands $B_j = \{h_j < y < h_{j-1}\}$ over $\sigma$ for the $j$ with $d_j > 0$, and of the graphs over $\sigma$ of the functions $h_j$ that lie strictly between $f$ and $g$. A graph of an affine function over $\sigma$ is an open simplex of dimension $k$. If $d_j > 0$ then $B_j$ is the open simplex of dimension $k+1$ with vertices
\[
p_i = (v_i, h_j(v_i)) \quad (0 \leq i \leq k), \qquad q = (v_j, g(v_j)).
\]
Indeed, the $p_i$ are affinely independent and $q - p_j = (0, d_j)$ is vertical. A point $\sum_i s_i p_i + s q$ with positive coefficients of sum $1$ is $(x, h_j(x) + s d_j)$, where $x \in \sigma$ has barycentric coordinates $t_i = s_i$ for $i \neq j$ and $t_j = s_j + s$; and $s$ ranges over $0 < s < t_j$, that is, $h_j(x) < y < h_j(x) + t_j d_j = h_{j-1}(x)$.
\end{proof}

\begin{lem}[Triangulation]\label{tri}
Let $g_1, \ldots, g_r$ be rational affine functions on $\Gamma^n$ and let $D \subseteq \Gamma^n$ be a bounded Boolean combination of the sets $g_i > 0$, $g_i = 0$. Then $D$ is a finite disjoint union of open simplices with rational vertices, on each of which every $g_i$ has constant sign.
\end{lem}
\begin{proof}
Induction on $n$; for $n = 0$ there is nothing to prove. Let $n \geq 1$, write the coordinates as $(x, y)$ with $x \in \Gamma^{n-1}$, and choose $N \in \N$ with $D \subseteq [-N, N]^n$. We may add $y - N$ and $y + N$ to the list. Write $g_i = c_i y + a_i(x)$. For $c_i \neq 0$ put $\varphi_i = - a_i / c_i$, so that the sign of $g_i$ is determined by the sign of $y - \varphi_i(x)$. Apply the induction hypothesis to the box $[-N, N]^{n-1}$ and the functions $a_i$ with $c_i = 0$, the differences $\varphi_i - \varphi_{i'}$, and $x_j \pm N$. On each simplex $\sigma$ so obtained the functions $\varphi_i$ are ordered: their distinct restrictions are $\psi_1 < \cdots < \psi_s$ on $\sigma$, with $\psi_1 \leq -N$ and $\psi_s \geq N$. The cylinder over $\sigma$ is the disjoint union of the graphs of the $\psi_j$, the bands between consecutive $\psi_j$, and two rays, and every $g_i$ has constant sign on each of these. So $D$ is a union of some of them; the rays do not occur, since they lie outside $[-N, N]^n$. The graphs are open simplices, and the bands are disjoint unions of open simplices by Lemma~\ref{prism}.
\end{proof}

\begin{lem}[One simplex]\label{cone}
Let $\sigma \subseteq \Gamma^n$ be an open simplex of dimension $k$ with vertices $w_0, \ldots, w_k \in \Q^n$ and let $\ell(x) = a \cdot x + b$ on $\sigma$ with $a \in \Z^n$, $b \in \Z$. Then $Z(\sigma, \ell) \in R[T]_\dag$, $\deg Z(\sigma, \ell) = 0$ and $\lim_{T \to \infty} Z(\sigma, \ell) = (-1)^{k+1}$.
\end{lem}
\begin{proof}
Put $L(x, m) = a \cdot x + b m$ on $\Z^{n+1}$. The map $\gamma \mapsto (m\gamma, m)$ identifies the index set of the double sum with
\[
C = \{(x, m) \in \Z^n \times \Z : m \geq 1, \ x/m \in \sigma\}, \qquad Z(\sigma, \ell) = \sum_{(x,m) \in C} U^{-L(x,m)} T^m .
\]
Choose integers $m_j \geq 1$ with $m_j w_j \in \Z^n$ and put $v_j = (m_j w_j, m_j) \in \Z^{n+1}$. The $v_j$ are linearly independent, and $C$ is the set of lattice points $\sum_j \lambda_j v_j$ with all $\lambda_j > 0$: if $x/m = \sum_j t_j w_j$ then $(x, m) = \sum_j (m t_j / m_j) v_j$, and conversely $m = \sum_j \lambda_j m_j$ is a positive integer.

Each $\lambda_j > 0$ is uniquely $n_j + \mu_j$ with $n_j \in \N$ and $0 < \mu_j \leq 1$. So every point of $C$ is uniquely $p + \sum_j n_j v_j$ with $n_j \in \N$ and $p$ in the finite set
\[
\Pi = \Bigl\{ \sum\nolimits_j \mu_j v_j : 0 < \mu_j \leq 1 \Bigr\} \cap \Z^{n+1},
\]
and every such sum is in $C$. Since $L$ and $m$ are additive, summing the geometric series gives
\[
Z(\sigma, \ell) = \frac{\sum_{p \in \Pi} U^{-L(p)}\, T^{m(p)}}{\prod_{j=0}^{k} \bigl(1 - U^{-L(v_j)}\, T^{m_j}\bigr)} .
\]
The denominator has degree $M = \sum_j m_j$ and leading coefficient $(-1)^{k+1} U^{-L(\sum_j v_j)}$. For $p = \sum_j \mu_j v_j \in \Pi$ we have $m(p) = \sum_j \mu_j m_j \leq M$, with equality only for $p = \sum_j v_j$. So the numerator has degree $M$ and leading coefficient $U^{-L(\sum_j v_j)}$.
\end{proof}

\begin{proof}[Proof of Lemma 8.5.2]
Apply Lemma~\ref{tri} to $\Delta$ and to the affine functions defining $\Delta$ and the pieces of the partition for $\ell$. We get $\Delta = \bigsqcup_\sigma \sigma$ with $\ell$ integral affine on each $\sigma$. Then $Z(\Delta, \ell) = \sum_\sigma Z(\sigma, \ell)$, and by Lemma~\ref{cone} and additivity of $\lim$ and of $\chi$,
\[
\lim_{T \to \infty} Z(\Delta, \ell) = \sum_\sigma (-1)^{\dim \sigma + 1} = - \sum_\sigma \chi(\sigma) = - \chi(\Delta). \qedhere
\]
\end{proof}

\begin{rem}
The Hadamard product (Lemma 8.4.1 of Hrushovski--Loeser) and the generation lemma of Hrushovski--Kazhdan are not used. Lemma~\ref{cone} is the classical generating function of the lattice points of a simplicial cone. Example: for $\Delta = [0,1)$ and $\ell = 0$ the simplices are $\{0\}$ and $(0,1)$, with series $T/(1-T)$ and $T^2/(1-T)^2$ and limits $-1$ and $1$; the sum is $0 = -\chi([0,1))$.
\end{rem}
\end{document}
